\documentclass[10pt,a4paper]{amsart}
\usepackage[margin=19mm]{geometry}
\usepackage{amsmath,amssymb,mathtools}
\usepackage[scr=rsfs,cal=euler]{mathalfa}
\usepackage{xfrac}
\usepackage[unicode,hidelinks,bookmarks=false]{hyperref}
\newcommand{\ie}{i.e.\ }
\newcommand{\cf}{cf.\ }
\newcommand{\resp}{respectively\ }
\newcommand{\Subsection}{\subsection}
\providecommand{\nobreakdash}{\nobreak\hbox{-}}
\setlength{\emergencystretch}{2em}
\multlinegap0pt
\usepackage{bm}
\usepackage{bbm}
\usepackage{dsfont}
\usepackage{xspace}
\usepackage{cancel}
\usepackage{mathtools}
\usepackage{amsthm}
\makeatletter
\@ifundefined{theorem}{\newtheorem{theorem}{Theorem}}{}
\makeatother
\title[Cohomologically Calibrated Affine Connections and Forced Irr]{Cohomologically Calibrated Affine Connections and Forced Irreducibility}
\author{Alexander Pigazzini, Magdalena Toda}
\date{Source: arXiv:2508.11024v2 (CC BY 4.0)}
\begin{document}
\maketitle

\noindent\emph{Source and licence.} Cohomologically Calibrated Affine Connections and Forced Irreducibility. Version arXiv:2508.11024v2 (CC BY 4.0), distributed under the Creative Commons Attribution 4.0 International licence, as recorded in the arXiv licence metadata for this exact version and shown on the abstract page https://arxiv.org/abs/2508.11024v2. Full rights notice: licenses/arxiv-2508.11024-CC-BY-4.0.txt.\par\smallskip

\noindent\textit{Editorial scope.} Counted unit: the Theorem whose source label is \texttt{thm:main} in the article (source0.tex lines 83--85, source label \texttt{thm:main}), delivered together with the article's own complete original proof of it (source0.tex lines 87--101, one \texttt{proof} environment of 15 lines). No mathematics is written, repaired or re-derived: statement and proof are the article's own text line for line. Required context retained uncounted: none. Every definition, display and label of those lines is retained unchanged. Omitted from this package and not redistributed: 1--82, 86--86, 102--228. Nothing inside a retained excerpt is omitted or altered: every retained body is delivered exactly as the article's lines are written, so no omission receipt of any kind accompanies this package. Citations: the retained excerpts carry no citation command at all, so no bibliography entry of the article is transcribed and no reference is redistributed. Reference adaptation: none was needed and none was made; every reference inside the delivered unit resolves inside it, so no unresolved reference or citation is produced. Printed slips kept literal: no printed word, symbol, hypothesis or number of the counted statement or of its proof was altered. Layout and class: the article's own class is \texttt{article} with packages including amsmath, amssymb, amsthm, mathrsfs, hyperref, geometry; the wrapper is the lane's pinned \texttt{amsart} editorial wrapper at 10pt on A4, which restates the theorem-like environments the retained text needs and adds the article's own macro definitions that the retained text needs (none), with extra wrapper packages (bm, bbm, dsfont, xspace, cancel, mathtools). The delivered numbering is the unit's own. Assets: no retained span contains a figure environment, a graphics-inclusion command, a PDF-page inclusion, a table or a TikZ picture; no image file is copied into this package, the package declares no asset dependency and no image file is redistributed with it. Licence: version arXiv:2508.11024v2 of this article is distributed under the Creative Commons Attribution 4.0 International licence, as recorded in the arXiv licence metadata for this exact version and shown on the abstract page https://arxiv.org/abs/2508.11024v2. A full rights notice is delivered at licenses/arxiv-2508.11024-CC-BY-4.0.txt. Characters: every retained excerpt is pure ASCII, so the delivered engine, XeLaTeX with its default Latin Modern text font, reports no missing character.
\begin{theorem} \label{thm:main}
Let $M = M_1 \times M_2$ be a product of compact, oriented Riemannian manifolds. Let $[\omega] \in H^3(M; \mathbb{R})$ be a non-trivial mixed cohomology class. Then, for any admissible connection $\nabla \in \mathcal{T}_\omega$, the holonomy algebra $\mathrm{hol}(\nabla)$ is irreducible with respect to the decomposition $V_1 \oplus V_2$.
\end{theorem}

\begin{proof}
The holonomy algebra $\mathrm{hol}(\nabla)$ is the Lie algebra generated by the curvature operators $R(X,Y)$. It is irreducible if and only if the curvature tensor $R(\nabla)$ does not preserve the subspaces $V_1$ and $V_2$. The curvature of $\nabla = \nabla^{\mathrm{LC}} + T$ is given by the standard identity:
\begin{equation} \label{eq:full_curvature_decomposition}
R(\nabla)(X,Y)Z = R^{\mathrm{LC}}(X,Y)Z + (\nabla^{\mathrm{LC}}_X T)(Y,Z) - (\nabla^{\mathrm{LC}}_Y T)(X,Z) + T(T(X,Y),Z) - T(T(X,Z),Y).
\end{equation}
The Levi-Civita term $R^{\mathrm{LC}}$ preserves the decomposition. We will show that the torsion-dependent terms, collectively denoted $R_T$, do not. The proof proceeds by demonstrating that for any admissible torsion $T \in \mathcal{T}_\omega$, the curvature tensor $R(\nabla)$ must have non-zero off-diagonal blocks. This is a structural consequence of the calibration, established via three steps:
\begin{enumerate}
    \item \textbf{Structure of Torsion:} A mixed class $[\omega]$ forces the harmonic part of the torsion, $T_h$, to have a specific algebraic structure that inextricably links $V_1$ and $V_2$. This structure is intrinsic to the cohomology class.

    \item \textbf{Structure of Curvature:} This "mixing" structure of $T_h$ generates a robustly non-zero off-diagonal component in the harmonic part of the curvature tensor, $R_{T_h}$. Specifically, curvature operators of the form $R(\nabla)(X,Y)$ with $X \in V_1, Y \in V_2$ act as non-zero maps between the subspaces $V_1$ and $V_2$.

    \item \textbf{Formal Non-Cancellation Argument:} The core of the proof is to show that the off-diagonal components generated by the harmonic part $T_h$ cannot be globally cancelled by the terms involving the non-harmonic part $T_{nh} = T - T_h$. We prove this formally by leveraging the $L^2$-orthogonality of the Hodge decomposition. A cancellation would imply that a topologically non-trivial harmonic tensor is equal to a topologically trivial one, which is impossible in the $L^2$ sense.
\end{enumerate}
A detailed, constructive realization of this proof for the case $M=S^2 \times \Sigma_g$ is provided in Section 4 and the Appendix. This explicit analysis demonstrates the mechanism in a concrete setting, confirming that the off-diagonal curvature components are robustly non-zero for any admissible connection. Since $R(\nabla)$ does not preserve the decomposition, the holonomy algebra $\mathrm{hol}(\nabla)$ is necessarily irreducible.
\end{proof}

\end{document}
